\documentclass[dvipsnames]{beamer}

\usepackage[T2A]{fontenc}
\usepackage[utf8]{inputenc}
\IfFileExists{mongolian.ldf}{%
  \usepackage[english,mongolian]{babel}%
}{%
  \usepackage[english]{babel}%
}

\usefonttheme[onlymath]{serif}
\usetheme{Warsaw}
\setbeamertemplate{navigation symbols}{}

\usepackage{amsmath}
\usepackage{amssymb}
\usepackage{graphicx}

\title[ARMA203]{ARMA203}
\subtitle{Бүлэг 4: Туршилтын загварчлал}
\author[А.~Мэнд-Оргил, М.~Номин]{А.~Мэнд-Оргил, М.~Номин}
\institute[МУИС]{Статистик, Хэрэглээний математикийн тэнхим\\Мэдээллийн технологи, электроникийн сургууль\\Монгол Улсын Их Сургууль}
\date{}

\newcommand{\safeimage}[2][]{%
  \IfFileExists{#2}{%
    \includegraphics[#1]{#2}%
  }{%
    \setlength{\fboxsep}{6pt}%
    \fbox{%
      \parbox[c][0.22\textheight][c]{0.78\linewidth}{%
        \centering
        \texttt{#2}\\
        Зураг олдсонгүй.
      }%
    }%
  }%
}

\begin{document}

%---------------------------------
\begin{frame}
\titlepage
\end{frame}

%---------------------------------
\section{Хичээлийн агуулга}

\begin{frame}{9-р хичээлийн агуулга}
\begin{enumerate}
\item Оршил
\item Chesapeake Bay булан дахь агнуур ба нэг гишүүнт загварууд
\item Цэнхэр загас агнах
\item Цэнхэр хавч агнах
\item Өндөр эрэмбийн олон гишүүнт загварууд
\item Дуу хураагуурын хугацаа
\end{enumerate}
\end{frame}

%---------------------------------
\section{Оршил}

\begin{frame}{Оршил}
\small
3-р бүлэгт бид муруй тааруулах (\emph{curve fitting}) болон
интерполяци (\emph{interpolation})-ийн философийн ялгааг авч үзсэн.
Муруй тааруулах үед загвар бүтээгч (\emph{modeler}) ажиглагдаж буй
зан төлөвийг тайлбарлах тодорхой загварын төрлийг сонгох зарим
таамаглал гаргадаг.

\medskip

Хэрэв цуглуулсан өгөгдөл эдгээр таамаглал үндэслэлтэй болохыг баталж
байвал, загвар бүтээгч өгөгдөлд хамгийн сайн тохирох муруйн
параметрүүдийг тодорхой шалгуурын дагуу (жишээлбэл, хамгийн бага
квадратын арга --- \emph{least squares}) сонгодог.
\end{frame}

%---------------------------------
\begin{frame}{Оршил}
\small
Үндсэндээ загвар бүтээгчийн зорилго нь тодорхой таамаглал дээр үндэслэн
загвар сонгохоос илүүтэй цуглуулсан өгөгдөлд тулгуурласан туршилтын
(\emph{эмпирик}) загвар байгуулах явдал юм.

\medskip

Ийм тохиолдолд загвар бүтээгч өгөгдөл цуглуулах болон шинжлэх үйл
явцад ихээхэн анхаардаг. Дараа нь өгөгдлийн цэгүүдийн хооронд таамаглал
хийхийн тулд өгөгдлийн ерөнхий чиг хандлагыг илэрхийлэх муруйг хайж
олдог.
\end{frame}

%---------------------------------
\begin{frame}{Оршил (үргэлжлэл)}
\small
Жишээлбэл, Зураг~4.1(a)-д өгөгдлийн цэгүүдийг авч үзье.
Хэрэв загвар бүтээгч квадрат (\emph{quadratic}) загвар гэж үзвэл,
өгөгдлийн цэгүүдэд парабола тааруулж болно. Үүнийг Зураг~4.1(b)-д
үзүүлсэн.

\medskip

Харин загвар бүтээгч тодорхой загварын төрлийг урьдчилан таамаглах
шалтгаангүй бол өгөгдлийн цэгүүдээр гөлгөр муруй дамжуулж болно.
Үүнийг Зураг~4.1(c)-д харуулсан.
\end{frame}

%---------------------------------
\begin{frame}{Зураг 4.1}
\begin{figure}
\centering
\safeimage[width=0.9\textwidth]{figure41.png}

\smallskip
\small
Хэрэв загвар бүтээгч квадрат хамаарал гэж таамаглавал өгөгдөлд
парабола тааруулж болно (b). Харин тодорхой таамаглал байхгүй бол
өгөгдлийн цэгүүдээр гөлгөр муруй дамжуулж болно (c).
\end{figure}
\end{frame}

%---------------------------------
\begin{frame}{4-р бүлэг: Туршилтын загварчлал}
\small
Энэ бүлэгт бид эмпирик (\emph{туршилтад суурилсан}) загвар байгуулах
асуудлыг авч үзнэ.

\medskip
4.1-р хэсэгт өгөгдлийн чиг хандлагыг барьж авах нэг гишүүнт
(\emph{one-term}) энгийн загвар сонгох процессыг судална.

\medskip
Зарим нөхцөлд загвар бүтээгчид ил тод математик загвар байгуулах
ороллого хийх нь ховор байдаг. Жишээлбэл, Chesapeake Bay дахь далайн
амьтдын агнуурын хэмжээг урьдчилан таамаглах зэрэг.
\end{frame}

%---------------------------------
\begin{frame}{4-р бүлгийн агуулга}
\begin{itemize}
\item \textbf{4.2} Өндөр зэрэгт олон гишүүнт (\emph{higher-order polynomial})
      загварууд
\item \textbf{4.3} Бага зэрэгтэй олон гишүүнтүүд ашиглан өгөгдлийг жигд
      болгох (\emph{smoothing})
\item \textbf{4.4} Кубик сплайн интерполяц (\emph{cubic spline interpolation})
\end{itemize}

\medskip
Эдгээр аргуудаар өгөгдлийн чиг хандлагыг илүү нарийн загварчлах
боломжтой.
\end{frame}

%---------------------------------
\section{4.1 Chesapeake Bay дахь агнуур}

\begin{frame}{4.1 Chesapeake Bay дахь агнуур ба нэг гишүүнт загварууд}
\small
\textbf{(Harvesting in the Chesapeake Bay and Other One-Term Models)}

\medskip
Загвар бүтээгч тодорхой хэмжээний өгөгдөл цуглуулсан боловч ил тод
математик загвар байгуулж чадахгүй байгаа нөхцөлийг авч үзье.

\medskip
1992 онд \emph{Daily Press} сонин Виржиниа мужид хэвлэгдсэн бөгөөд
Chesapeake Bay-д далайн амьтан агнах талаар сүүлийн 50 жилийн
ажиглалтын өгөгдлийг нийтэлсэн.
\end{frame}

%---------------------------------
\begin{frame}{Chesapeake Bay агнуурын өгөгдөл}
\small
Бид дараах хоёр төрлийн өгөгдлийг ашиглан хэд хэдэн нөхцөлийг судална.

\begin{itemize}
\item bluefish (\textbf{цэнхэр загас}) агнуур
\item blue crab (\textbf{цэнхэр хавч}) агнуур
\end{itemize}

Эдгээр нь Chesapeake Bay-ийн арилжааны загас агнуурын салбарын мэдээлэл
юм. Хүснэгт~4.1 нь нэг гишүүнт загвар байгуулахад ашиглах өгөгдлийг
харуулна.
\end{frame}

%---------------------------------
\begin{frame}{Өгөгдлийн чиг хандлага}
\small
Bluefish агнуурын тархалтын график \textbf{Зураг~4.2}-т, харин
blue crab агнуурын тархалтын график \textbf{Зураг~4.3}-т үзүүлсэн.

\begin{itemize}
\item Зураг~4.2-оос харахад цаг хугацаа өнгөрөх тусам bluefish агнуур
      нэмэгдэх хандлагатай.
\item Энэ нь bluefish-ийн нөөц нэмэгдэж байгааг илэрхийлж болох юм.
\item Харин blue crab агнуур мөн өсөх хандлагатай боловч математик
      загвараар илэрхийлэх нь амар биш.
\end{itemize}

Бидний стратеги бол Хүснэгт~4.1-ийн өгөгдлийг хувиргах бөгөөд ингэснээр
график нь шулуунд ойролцоо харагдана.
\end{frame}

%---------------------------------
\begin{frame}{Хүснэгт 4.1}
\begin{figure}
\centering
\safeimage[width=0.9\textwidth]{figure42.png}
\end{figure}
\end{frame}

%---------------------------------
\begin{frame}{Зураг 4.2: Bluefish агнуур}
\begin{figure}
\centering
\safeimage[width=0.8\textwidth]{figure422.png}

\smallskip
\small
1940--1990 оны хооронд (5 жилийн интервалтай) bluefish агнуурын хэмжээ
ба суурь жилийн хоорондох тархалтын график.
\end{figure}
\end{frame}

%---------------------------------
\begin{frame}{Зураг 4.3: Blue crab агнуур}
\begin{figure}
\centering
\safeimage[width=0.8\textwidth]{figure43.png}

\smallskip
\small
1940--1990 оны хооронд (5 жилийн интервалтай) blue crab агнуурын хэмжээ
ба суурь жилийн хоорондох тархалтын график.
\end{figure}
\end{frame}

%---------------------------------
\begin{frame}{Зэрэгтийн шат (Ladder of Powers)}
\small
\begin{columns}[T]
\column{0.35\textwidth}
\centering
\safeimage[width=\textwidth]{ladder.png}

\column{0.65\textwidth}
Бид хувиргалтыг (\emph{transformation}) хэрхэн сонгох вэ? Үүнийг
тодорхойлохын тулд хувьсагч $z$-ийн зэрэгтийн шат
(\emph{ladder of powers})-ыг ашиглаж, өгөгдлийг шугаман хэлбэрт
ойртуулах тохиромжтой хувиргалтыг сонгоно.

\medskip
Зураг~4.4-т таван өгөгдлийн цэг $(x, y)$ болон $y=x$ шулууг
($x>1$ үед) хамт үзүүлсэн байна.

\medskip
Хэрэв бид цэг бүрийн $y$ утгыг $\sqrt{y}$ болгон өөрчилбөл,
$y$-ийн утгууд хоорондоо илүү ойр болно. Энэ үйлдэл нь
$y=\sqrt{x}$ гэсэн шинэ хамаарлыг бий болгоно.
\end{columns}
\end{frame}

%---------------------------------
\begin{frame}{Зэрэгтийн шатны нөлөө}
\small
Энэ мужид бүх $y$ утга багасах бөгөөд том утгууд нь жижиг
утгуудаас илүү их хэмжээгээр багасна.

\medskip
Харин $y$-ийн утгыг $\log y$ болгон өөрчлөх нь $\sqrt{y}$-ээс ч
илүү хүчтэй нөлөөтэй байдаг.

\medskip
Иймээс зэрэгтийн шатанд нэг алхам нэмэгдэх тусам өмнөхөөсөө илүү
хүчтэй нөлөөтэй хувиргалт бий болдог гэж үздэг.
\end{frame}

%---------------------------------
\begin{frame}{Зураг 4.4: Хувиргалтын харьцангуй нөлөө}
\begin{figure}
\centering
\safeimage[width=0.8\textwidth]{figure44.png}

\smallskip
\small
Энэ зураг нь анхны $y=x$ шулуун функц дээр дараах хувиргалтууд
хийхэд график хэрхэн өөрчлөгдөхийг харуулж байна.
\begin{itemize}
\item $y$-ийн утгыг $\sqrt{y}$ болгон солих
\item $y$-ийн утгыг $\log y$ болгон солих
\item $y$-ийн утгыг $\frac{1}{y}$ болгон солих
\end{itemize}
\end{figure}
\end{frame}

%---------------------------------
\begin{frame}{Хувиргалтын аргууд}
\small
Хэрэв бид доошоо хотойсон (\emph{concave down}), эерэг утгатай өсөх
функц авбал, жишээлбэл
\[
y=f(x), \quad x>1,
\]
графикийг илүү шугаман болгохын тулд дараах аргуудыг хэрэглэж болно.

\begin{enumerate}
\item \textbf{Баруун талын сүүлийг дээш сунгах.}

      Үүний тулд $y$ утгуудыг
      \[
      y^2, \; y^3
      \]
      гэх мэт болгон өөрчилж болно.

\item \textbf{Баруун талын сүүлийг зүүн тийш шахах.}

      Үүний тулд $x$ утгуудыг
      \[
      \sqrt{x}, \; \log x
      \]
      гэх мэт болгон өөрчилж болно.
\end{enumerate}
\end{frame}

%---------------------------------
\begin{frame}{Хувиргалтын анхаарах зүйл}
\small
Гэхдээ $\frac{1}{z}$ эсвэл $\frac{1}{z^2}$ гэх мэт хувиргалт хийхэд
нэг сул тал бий. Учир нь:

\begin{itemize}
\item өсөх функц буурах функц болж хувирч болно.
\end{itemize}

Иймээс өгөгдлийн шинжилгээнд ихэвчлэн $\log z$-оос доош байрлах
хувиргалтуудыг хэрэглэхдээ \textbf{сөрөг тэмдэг} ашигладаг.

\medskip
Ингэснээр шинэ өгөгдөл анхны өгөгдлийн адил өсөх дараалалтай хэвээр
үлддэг.

\medskip
Хүснэгт~4.2 нь практикт хамгийн их хэрэглэгддэг
\textbf{хувиргалтын шатлал}-ыг харуулж байна.
\end{frame}

%---------------------------------
\begin{frame}{Жишээ 1: Bluefish агнуур}
\small
Зураг~4.2-ын тархалтын графикаас харахад өгөгдлийн ерөнхий чиг
хандлага өсөх бөгөөд \emph{concave up} байна.

\medskip
Иймээс \textbf{зэрэгтийн шат} (\emph{ladder of powers}) ашиглан баруун
талын сүүлийг доош шахах хэрэгтэй.

\medskip
Үүний нэг арга нь
\[
y \rightarrow \log y
\]
хувиргалтыг хийх явдал юм.

\medskip
Өөр нэг боломж бол $x$ утгуудыг
\[
x^2, \; x^3
\]
гэх мэт өндөр зэрэгт хувиргах арга юм.
\end{frame}

%---------------------------------
\begin{frame}{Өгөгдлийн хувиргалт}
\small
Зураг~4.2 дээрх өгөгдлийн чиг хандлагыг эргэн санацгаая. Зэрэгтийн шат
(\emph{ladder of powers})-ыг ашиглан баруун гар талыг доош нь шахахын
тулд $y$-г $\log y$-ээр солих эсвэл шатаар доошилсон бусад
хувиргалтуудыг хэрэглэж болно.

\medskip
Өөр нэг сонголт бол $x$ утгуудыг
\[
x^2, \quad x^3
\]
гэх мэт өндөр зэрэгт хувиргах явдал юм.

\medskip
Бид Хүснэгт~4.3-т байгаа өгөгдлийг ашиглана. Энд 1940 оныг суурь жил
гэж авч $x=0$ гэж тэмдэглэсэн. Дараагийн жилүүд нь 5 жилийн зайтай
нэмэгдэнэ.
\end{frame}

%---------------------------------
\begin{frame}{Хувиргалтын сонголт}
\small
Бид шатаар өгсөх ($x^2, x^3$ гэх мэт) янз бүрийн утгууд руу $x$-г
өөрчлөх замаар баруун гар талын сүүлийг баруун тийш шахаж эхэлдэг.

\medskip
Гэвч эдгээр хувиргалтын аль нь ч \textbf{шугаман график}
үүсгэдэггүй.

\medskip
Дараа нь бид $y$-г
\[
\sqrt{y}, \quad \log y
\]
хувиргалтуудаар сольж, шатаар доошилно.

\medskip
$\sqrt{y}$ болон $\log y$ графикуудын аль аль нь $x$-тэй
харьцуулахад илүү \textbf{шугаман} харагдаж байв. Иймээс бид дараах
загварыг сонгоно:
\[
\log y \;\text{vs}\; x
\]
\end{frame}

%---------------------------------
\begin{frame}{Хамгийн бага квадратын арга}
\small
Бид дараах хэлбэрийн загварыг тааруулсан:
\[
\log y = mx + b
\]

Хамгийн бага квадратын аргыг ашигласны үр дүнд дараах ойролцоолсон
тэгшитгэл гарсан:
\[
\log y = 0.7231 + 0.1654x
\]

\begin{itemize}
\item $x$ --- суурь жил
\item $y$ --- bluefish агнуурын хэмжээ
\item $y$-г $10^{4}$ фунтээр хэмжсэн
\end{itemize}

Үүнийг экспоненциал хэлбэрт шилжүүлбэл
\[
y = 5.2857(1.4635)^x
\]
\end{frame}

%---------------------------------
\begin{frame}{Зураг 4.5: Bluefish загвар}
\begin{figure}
\centering
\safeimage[width=0.8\textwidth]{figure45.png}

\smallskip
\small
Энд $y$ нь bluefish агнуурын хэмжээг $10^{4}$ фунтээр хэмжсэн утга,
харин $x$ нь суурь жил юм. Муруйн график нь өгөгдөлд боломжийн сайн
тохирч байна.
\end{figure}
\end{frame}

%---------------------------------
\begin{frame}{Жишээ 2: Blue crab агнуур}
\small
Өмнө дурдсанчлан, бид өгөгдлийг шугаман хэлбэрт ойртуулахын тулд
$y$-ийн утгуудыг $y^2$, $y^3$ гэх мэт утгуудаар сольж, эсвэл зэрэгтийн
шат (\emph{ladder of powers})-аар дээшилсэн бусад хувиргалтуудыг
туршиж үзэж болно.

\medskip
Хэд хэдэн туршилтын дараа бид $x$-ийн утгуудыг $\sqrt{x}$ утгаар
солихыг сонгосон. Энэ хувиргалт нь графикийн баруун талын сүүлийг
зүүн тийш шахдаг.

\medskip
Зураг~4.6-д $y$-г $\sqrt{x}$-тэй харьцуулсан графикийг үзүүлсэн.
\end{frame}

%---------------------------------
\begin{frame}{Blue crab загвар}
\small
Эхлэл цэгээр дамжин өнгөрөх шулуу:
\[
y = k\sqrt{x}
\]
хамгийн бага квадратын аргаар $k$-ийн утгыг олбол
\[
y = 158.344\sqrt{x}
\]

Энд
\begin{itemize}
\item $y$ --- blue crab агнуур ($10^{4}$ фунт)
\item $x$ --- суурь жил
\end{itemize}

\vspace{0.3cm}
\centering
\safeimage[width=0.6\textwidth]{table44.png}

\smallskip
\small Хүснэгт~4.4: Blue crab өгөгдөл
\end{frame}

%---------------------------------
\begin{frame}{Blue crab өгөгдөл}
\centering
\safeimage[width=\textwidth,height=0.6\textheight,keepaspectratio]{table44.png}
\end{frame}

%---------------------------------
\begin{frame}{Зураг 4.6: Blue crab ба $\sqrt{x}$}
\centering
\safeimage[width=0.9\textwidth]{figure46.png}

\smallskip
\small
Blue crab агнуур ($10^{4}$ lb) ба $\sqrt{x}$-ийн хоорондох тархалтын
график.
\end{frame}

%---------------------------------
\begin{frame}{Зураг 4.7: $y = 158.344\sqrt{x}$}
\centering
\safeimage[width=0.9\textwidth]{figure47.png}

\smallskip
\small
Эхлэл цэгээр дайрсан $y = k\sqrt{x}$ шулуу. Хамгийн бага квадратын
аргаар $k = 158.344$ гарсан.
\end{frame}

%---------------------------------
\begin{frame}{Зураг 4.8: Загвар ба өгөгдөл}
\centering
\safeimage[width=0.9\textwidth]{figure48.png}

\smallskip
\small
Өгөгдлийн тархалтын график дээр $y = 158.344\sqrt{x}$ загварын муруйг
давхардуулан харуулсан.
\end{frame}

%---------------------------------
\begin{frame}{Загваруудыг шалгах}
\small
Загвар дээр үндэслэн хийсэн таамаглал хэр сайн бэ? Үүний нэг хариулт
нь ажиглагдсан утгуудыг таамагласан утгуудтай харьцуулах явдал юм.

\medskip
Бид өгөгдлийн хос бүрийн хувьд үлдэгдэл (\emph{residual}) болон
харьцангуй алдааг тооцож болно.

\medskip
Ихэнх тохиолдолд загвар ашиглан ирээдүйн утгуудыг таамаглах эсвэл
экстраполяци хийх шаардлага гардаг. Жишээлбэл, 2010 онд Chesapeake Bay
дахь агнуурын хэмжээг эдгээр загварууд хэр сайн таамаглах вэ?

\[
\text{Bluefish } y = 5.2857(1.4635)^{14}
= 1092.95 \; (10^4 \text{ lb})
\approx 10.9 \text{ million lb}
\]

\[
\text{Blue crabs } y = 158.344\sqrt{14}
= 592.469 \; (10^4 \text{ lb})
\approx 5.92 \text{ million lb}
\]
\end{frame}

%---------------------------------
\begin{frame}{Загварын үнэлгээ}
\small
Эдгээр энгийн нэг гишүүнт загваруудыг ихэвчлэн \textbf{интерполяцид}
ашиглах нь тохиромжтой, харин \textbf{экстраполяцид} ашиглах нь
төдийлөн найдвартай биш байж болно.

\medskip
Эмпирик загвар байгуулахдаа бид хамгийн түрүүнд цуглуулсан өгөгдлийг
анхааралтай шинжилдэг.

\begin{itemize}
\item Өгөгдөлд тодорхой чиг хандлага байна уу?
\item Чиг хандлагаас хэт зөрсөн (\emph{outlier}) цэгүүд байна уу?
\item Хэрэв ийм цэгүүд байвал өгөгдлийн алдаа эсэхийг шалгах хэрэгтэй.
\end{itemize}

Хэрэв чиг хандлага тодорхой байвал өгөгдлийг шугаман хэлбэрт
ойртуулах тохиромжтой хувиргалтыг сонгоно.
\end{frame}

%---------------------------------
\section{4.2 Өндөр эрэмбийн олон гишүүнт загварууд}

\begin{frame}{4.2 Өндөр эрэмбийн олон гишүүнт загварууд}
\small
4.1-р хэсэгт бид өгөгдлийн чиг хандлагыг тайлбарлах энгийн
\textbf{нэг гишүүнт загвар} олох боломжийг судалсан.

\medskip
Ийм загварууд нь энгийн бүтэцтэй тул
\begin{itemize}
\item загварын мэдрэмжийн шинжилгээ,
\item оновчлол,
\item өөрчлөлтийн хурд,
\item муруйн доорх талбай
\end{itemize}
зэрэг анализ хийхэд тохиромжтой.

\medskip
Гэвч нэг гишүүнт загварууд нь өгөгдлийн бүх чиг хандлагыг тайлбарлах
боломжгүй байдаг. Тиймээс зарим тохиолдолд \textbf{олон гишүүнт
загвар} ашиглах шаардлага гардаг.
\end{frame}

%---------------------------------
\begin{frame}{Олон гишүүнт загвар}
\small
Энэ хэсэгт бид \textbf{олон гишүүнт загвар}-ын нэг төрлийг судална.

\medskip
Олон гишүүнт функцууд нь
\begin{itemize}
\item интеграл авахад хялбар,
\item уламжлал олоход хялбар.
\end{itemize}

\medskip
Гэвч тэдгээр нь сул талтай. Жишээлбэл, босоо асимптоттой өгөгдлийг нэг
олон гишүүнтээр ойролцоо дүрслэх нь тохиромжгүй байдаг. Ийм тохиолдолд
олон гишүүнтийн харьцаа
\[
\frac{p(x)}{q(x)}
\]
ашиглах нь илүү тохиромжтой.
\end{frame}

%---------------------------------
\begin{frame}{Олон гишүүнт ба өгөгдлийн цэгүүд}
\small
Одоо дараах асуудлыг авч үзье.

\medskip
Хоёр өгөгдлийн цэг өгөгдсөн гэж бодъё:
\[
(x_1,y_1), \quad (x_2,y_2)
\]

Эдгээр хоёр цэгээр дайрсан \textbf{цорын ганц шулуун} оршино.
\[
y = a_0 + a_1 x
\]

Шулуу нь дараах нөхцөлийг хангана:
\[
y_1 = a_0 + a_1 x_1
\]
\[
y_2 = a_0 + a_1 x_2
\]

Үүнтэй адил \textbf{3 цэгээр} дайрсан \textbf{2-р зэрэгт олон
гишүүнт} байдаг.
\end{frame}

%---------------------------------
\begin{frame}{Зураг 4.10: Олон гишүүнт функц}
\centering
\safeimage[width=0.9\textwidth]{figure410.png}

\smallskip
\small
3 цэгээр 2-р зэрэгт олон гишүүнт дамжин өнгөрч болно.
\end{frame}

%---------------------------------
\begin{frame}{2-р зэрэгт олон гишүүнт}
\small
Үүнтэй адил гурван өөр өгөгдлийн цэгээр (Зураг~4.10(b)-д үзүүлсэн шиг)
хамгийн ихдээ 2-р зэрэгтэй \textbf{цорын ганц олон гишүүнт функц}
дамжин өнгөрч чадна.
\[
y = a_0 + a_1 x + a_2 x^2
\]

Коэффициентүүд $a_0, a_1, a_2$-ийг дараах тэгшитгэлийн системийг
шийдвэрлэж олно:
\[
y_1 = a_0 + a_1 x_1 + a_2 x_1^2
\]
\[
y_2 = a_0 + a_1 x_2 + a_2 x_2^2
\]
\[
y_3 = a_0 + a_1 x_3 + a_2 x_3^2
\]
\end{frame}

%---------------------------------
\begin{frame}{Олон гишүүнтийн онцлог}
\small
``\textbf{хамгийн ихдээ}'' (\emph{at most}) гэсэн тодорхойлолтыг
яагаад хэрэглэснийг тайлбарлая.

\medskip
Хэрэв гурван цэг нэг шулуун дээр байрлавал 2-р зэрэгт олон гишүүнт
биш, харин
\[
\text{1-р зэрэгт олон гишүүнт}
\]
өөрөөр хэлбэл шулуун гарна.

\medskip
Мөн ``\textbf{цорын ганц}'' гэсэн тодорхойлолт чухал. Учир нь 2-оос
их зэрэгтэй \textbf{хязгааргүй олон} олон гишүүнт эдгээр гурван
цэгээр дамжин өнгөрч чадна.
\end{frame}

%---------------------------------
\begin{frame}{Жишээ: Дуу хураагуурын хугацаа}
\small
Бид тодорхой дуу хураагуурын тоолуурын заалт болон тоглох хугацааны
өгөгдлийг цуглуулсан.

\begin{center}
\begin{tabular}{c c c c c c c c c}
$c_i$ & 100 & 200 & 300 & 400 & 500 & 600 & 700 & 800 \\
$t_i$ (sec) & 205 & 430 & 677 & 945 & 1233 & 1542 & 1872 & 2224
\end{tabular}
\end{center}

Энд
\begin{itemize}
\item $c_i$ --- тоолуурын заалт
\item $t_i$ --- тоглосон хугацаа (секунд)
\end{itemize}

Эдгээр өгөгдлөөр \textbf{эмпирик загвар} байгуулж болно.
\end{frame}

%---------------------------------
\begin{frame}{Олон гишүүнт загвар}
\small
8 өгөгдлийн цэг байгаа тул хамгийн ихдээ \textbf{7-р зэрэгт олон
гишүүнт} үүснэ.

\[
P_7(c)
= a_0 + a_1 c + a_2 c^2 + a_3 c^3
+ a_4 c^4 + a_5 c^5 + a_6 c^6 + a_7 c^7
\]

Коэффициентүүдийг шугаман тэгшитгэлийн системийг бодож олно.
\end{frame}

%---------------------------------
\begin{frame}{Олон гишүүнтийн коэффициентүүд}
\small
8 өгөгдлийн цэг байгаа тул коэффициентүүд $a_i$ дараах шугаман
тэгшитгэлийн системийг хангана:

\[
205 = a_0 + 1a_1 + 1^2 a_2 + 1^3 a_3 + 1^4 a_4 + 1^5 a_5 + 1^6 a_6 + 1^7 a_7
\]
\[
430 = a_0 + 2a_1 + 2^2 a_2 + 2^3 a_3 + 2^4 a_4 + 2^5 a_5 + 2^6 a_6 + 2^7 a_7
\]
\[
\vdots
\]
\[
2224 = a_0 + 8a_1 + 8^2 a_2 + 8^3 a_3 + 8^4 a_4 + 8^5 a_5 + 8^6 a_6 + 8^7 a_7
\]
\end{frame}

%---------------------------------
\begin{frame}{Тоон тооцооллын хүндрэл}
\small
Ийм том шугаман тэгшитгэлийн системийг өндөр нарийвчлалтай бодох нь
заримдаа хүндрэлтэй байдаг.

\begin{itemize}
\item Өгөгдлийг 7-р зэрэгт дэвшүүлэхэд маш их ялгаатай хэмжээнүүд үүсдэг.
\item Иймээс коэффициентүүдийг аль болох өндөр нарийвчлалтай тооцоолох
      шаардлагатай.
\item Жижиг коэффициент $x^7$ зэрэгтэй үржигдэхэд том утга болж хувирч
      болно.
\end{itemize}

Иймээс өгөгдлийн хэмжээг багасгахын тулд тоолуурын утгуудыг 100-д
хувааж ашигласан.
\end{frame}

%---------------------------------
\begin{frame}{Polynomial загварын коэффициентүүд}
\small
Тооцооллын үр дүнд дараах коэффициентүүд гарсан:
\[
a_0 = -13.9999923
\]
\[
a_1 = 232.9119031
\]
\[
a_2 = -29.08333188
\]
\[
a_3 = 19.78472156
\]
\[
a_4 = -5.354166491
\]
\[
a_5 = 0.8013888621
\]
\[
a_6 = -0.0624999978
\]
\[
a_7 = 0.0019841269
\]
\end{frame}

%---------------------------------
\begin{frame}{Олон гишүүнт загварын шалгалт}
\small
Одоо эмпирик загвар өгөгдөлд хэр сайн тохирч байгааг шалгаж үзье.
Полиномын таамагласан утгуудыг $P_7(c_i)$ гэж тэмдэглэвэл:

\begin{center}
\small
\begin{tabular}{c|cccccccc}
$c_i$ & 100 & 200 & 300 & 400 & 500 & 600 & 700 & 800 \\
\hline
$t_i$ & 205 & 430 & 677 & 945 & 1233 & 1542 & 1872 & 2224 \\
$P_7(c_i)$ & 205 & 430 & 677 & 945 & 1233 & 1542 & 1872 & 2224
\end{tabular}
\end{center}

$P_7(c_i)$-ийн утгуудыг 4 орны нарийвчлалтай тоймловол ажиглагдсан
утгуудтай бүрэн тохирч байна. Иймээс абсолют хазайлт $0$ гарна.
\end{frame}

%---------------------------------
\begin{frame}{Загварын үнэлгээ}
\small
Гэхдээ 3-р бүлэгт үзсэн \textbf{тааруулалтын шалгууруудыг} дангаар нь
ашиглах нь загварын чанарыг бүрэн дүгнэхэд хангалтгүй байж болно.

\medskip
Бид дараах асуултыг тавих хэрэгтэй.
\begin{itemize}
\item Энэ загвар бусад загвараас илүү сайн уу?
\item Өгөгдлийн чиг хандлагыг үнэхээр зөв илэрхийлж байна уу?
\end{itemize}

\medskip
Иймээс $P_7(c_i)$ загвар өгөгдлийн ерөнхий чиг хандлагыг хэрхэн барьж
байгааг графикаар харъя.
\end{frame}

%---------------------------------
\begin{frame}{Зураг 4.11: Tape recorder загвар}
\begin{figure}
\centering
\safeimage[width=0.8\textwidth]{figure411.png}
\end{figure}

\smallskip
\small
Дуу хураагуурын хугацааг урьдчилан таамаглах эмпирик олон гишүүнт
загвар.
\end{frame}

%---------------------------------
\begin{frame}{Lagrange-ын куб олон гишүүнт}
\small
Дараах өгөгдөл өгөгдсөн гэж үзье.

\begin{center}
\begin{tabular}{c c c c c}
$x$ & $x_1$ & $x_2$ & $x_3$ & $x_4$ \\
$y$ & $y_1$ & $y_2$ & $y_3$ & $y_4$
\end{tabular}
\end{center}

Куб олон гишүүнтийг дараах байдлаар бичиж болно:

\[
\resizebox{0.95\textwidth}{!}{$
P_3(x)=
\frac{(x-x_2)(x-x_3)(x-x_4)}{(x_1-x_2)(x_1-x_3)(x_1-x_4)}y_1
+
\frac{(x-x_1)(x-x_3)(x-x_4)}{(x_2-x_1)(x_2-x_3)(x_2-x_4)}y_2
+
\frac{(x-x_1)(x-x_2)(x-x_4)}{(x_3-x_1)(x_3-x_2)(x_3-x_4)}y_3
+
\frac{(x-x_1)(x-x_2)(x-x_3)}{(x_4-x_1)(x_4-x_2)(x_4-x_3)}y_4
$}
\]
\end{frame}

%---------------------------------
\begin{frame}{Lagrange-ын олон гишүүнтийн тайлбар}
\small
Энэ олон гишүүнт нь дараах шинжтэй:
\begin{itemize}
\item $x = x_i$ үед
      \[
      P_3(x_i)=y_i
      \]
\item Иймээс олон гишүүнт өгөгдлийн бүх цэгээр дамжина.
\item $x_i$ утгууд өөр өөр байх ёстой.
\item Хэрэв ижил байвал хуваарьт 0 үүснэ.
\end{itemize}

Энэ хэлбэрийг \textbf{Lagrangian form of polynomial} гэж нэрлэдэг.
\end{frame}

%---------------------------------
\begin{frame}{Теорем 1: Lagrange-ын олон гишүүнт}
\small
\begin{block}{Теорем}
Хэрэв
\[
(x_0,y_0),(x_1,y_1),\dots,(x_n,y_n)
\]
нь $(n+1)$ өөр өгөгдлийн цэгүүд бол хамгийн ихдээ $n$ зэрэгтэй
\textbf{цорын ганц олон гишүүнт} $P(x)$ оршино.

\[
P(x_k)=y_k, \quad k=0,1,\dots,n
\]

Энэ олон гишүүнт дараах хэлбэртэй:
\[
P(x)=y_0L_0(x)+y_1L_1(x)+\dots+y_nL_n(x)
\]
энд
\[
L_k(x)=
\frac{(x-x_0)(x-x_1)\cdots(x-x_{k-1})(x-x_{k+1})\cdots(x-x_n)}
{(x_k-x_0)(x_k-x_1)\cdots(x_k-x_{k-1})(x_k-x_{k+1})\cdots(x_k-x_n)}
\]
\end{block}
\end{frame}

%---------------------------------
\begin{frame}{Өндөр эрэмбийн олон гишүүнтийн давуу ба сул тал}
\small
Олон гишүүнт функц нь интеграл болон уламжлал авахад хялбар байдаг.

\medskip
Хэрэв олон гишүүнт функц нь өгөгдлийн үндсэн зан төлөвийг сайн
илэрхийлж байвал
\begin{itemize}
\item муруйн доорх талбайг ойролцоолох,
\item өөрчлөлтийн хурдыг тодорхойлох
\end{itemize}
боломжтой.

\medskip
Гэвч өндөр зэрэгт олон гишүүнтүүд нь зарим тохиолдолд сул талтай
байдаг. Жишээлбэл, Хүснэгт~4.10-д өгөгдсөн 17 өгөгдлийн цэгийн хувьд
\[
y = 0 \quad (-8 \le x \le 8)
\]
байна. Энэ нь өндөр эрэмбийн олон гишүүнт ашиглахад гарч болох
асуудлыг харуулдаг.
\end{frame}

%---------------------------------
\begin{frame}{Хүснэгт 4.10}
\centering
\small
\resizebox{\textwidth}{!}{%
\begin{tabular}{c|rrrrrrrrrrrrrrrrr}
$x_i$ & -8 & -7 & -6 & -5 & -4 & -3 & -2 & -1 & 0 & 1 & 2 & 3 & 4 & 5 & 6 & 7 & 8 \\
\hline
$y_i$ & 0 & 0 & 0 & 0 & 0 & 0 & 0 & 0 & 0 & 0 & 0 & 0 & 0 & 0 & 0 & 0 & 0
\end{tabular}%
}
\end{frame}

%---------------------------------
\begin{frame}{Өндөр эрэмбийн олон гишүүнтийн сул тал}
\centering
\safeimage[width=0.8\textwidth]{figure412.png}

\vspace{0.3cm}

\small
\textbf{Зураг 4.12.} Өгөгдлийн цэгүүдээр дайруулан өндөр эрэмбийн олон
гишүүнт тааруулахад хязгаарын орчимд хүчтэй хэлбэлзэл
(\emph{oscillation}) үүсдэг.
\end{frame}

%---------------------------------
\begin{frame}{Олон гишүүнтийн хэлбэлзэл}
\small
Өндөр эрэмбийн олон гишүүнт нь өгөгдлийн цэгүүдээр яг дайрч өнгөрдөг
боловч:
\begin{itemize}
\item интервалын төгсгөлүүдэд хүчтэй хэлбэлзэл үүсдэг;
\item $x=\pm 8$ орчимд алдаа ихэсдэг;
\item уламжлал болон талбай тооцоход алдаа гарна.
\end{itemize}

\vspace{0.3cm}
\textbf{Жишээ өгөгдөл:}
\[
\begin{array}{c|ccccccc}
x & 0.55 & 1.2 & 2 & 4 & 6.5 & 12 & 16 \\
\hline
y & 0.13 & 0.64 & 5.8 & 102 & 210 & 2030 & 3900
\end{array}
\]
\end{frame}

%---------------------------------
\begin{frame}{Зураг 4.13: Өгөгдлийн тархалт}
\centering
\safeimage[width=0.8\textwidth]{figure413.png}

\vspace{0.3cm}

\small
\textbf{Зураг 4.13.} Өгөгдлийн цэгүүдийн тархалт. Муруй нь өсөх бөгөөд
дээш хотойсон (\emph{concave up}) хэлбэртэй байна.
\end{frame}

%---------------------------------
\begin{frame}{Өндөр эрэмбийн олон гишүүнтийн сул тал}
\footnotesize
Бид өгөгдөлд 6-р зэрэгтэй олон гишүүнт тааруулна ($n=7$):
\[
\begin{aligned}
y ={}& -0.0138x^6 + 0.5084x^5 - 6.4279x^4 \\
     & + 34.8575x^3 - 73.9916x^2 + 64.3128x - 18.0951
\end{aligned}
\]

\vspace{0.3cm}
\begin{itemize}
\item Өндөр зэрэгт олон гишүүнт нь өгөгдөлд маш сайн таарна.
\item Гэвч интервалын төгсгөлд хэлбэлзэл үүсдэг.
\item Урьдчилан таамаглахад найдвартай биш.
\item Интерполяц ч мөн асуудалтай болдог.
\end{itemize}
\end{frame}

%---------------------------------
\begin{frame}{Зураг 4.14: Polynomial fit}
\centering
\safeimage[width=0.8\textwidth]{figure414.png}

\vspace{0.3cm}

\small
\textbf{Зураг 4.14.} 6-р зэрэгтэй олон гишүүнт өгөгдөл дээр
тааруулсан график.
\end{frame}

%---------------------------------
\begin{frame}{Хүснэгт 4.11: Өгөгдлийн багцууд}
\centering
\small
\begin{tabular}{c|ccccc}
$x_i$ & 0.2 & 0.3 & 0.4 & 0.6 & 0.9 \\
\hline
Case 1 & 2.7536 & 3.2411 & 3.8016 & 5.1536 & 7.8671 \\
Case 2 & 2.7540 & 3.2410 & 3.8020 & 5.1540 & 7.8670 \\
Case 3 & 2.7536 & 3.2411 & 3.8916 & 5.1536 & 7.8671
\end{tabular}

\vspace{0.3cm}
\footnotesize
Жижиг өөрчлөлт (алдаа) нь өгөгдөлд орсон байна.
\end{frame}

%---------------------------------
\begin{frame}{Олон гишүүнтийн коэффициентүүд}
\small
Дөрөвдүгээр зэрэгтэй олон гишүүнт:
\[
P_4(x)=a_0 + a_1 x + a_2 x^2 + a_3 x^3 + a_4 x^4
\]

\vspace{0.3cm}
\centering
\begin{tabular}{c|ccccc}
 & $a_0$ & $a_1$ & $a_2$ & $a_3$ & $a_4$ \\
\hline
Case 1 & 2 & 3 & 4 & -1 & 1 \\
Case 2 & 2.0123 & 2.8781 & 4.4159 & -1.5714 & 1.2698 \\
Case 3 & 3.4580 & -13.2000 & 64.7500 & -91.0000 & 46.0000
\end{tabular}

\vspace{0.3cm}
\footnotesize
Өгөгдлийн бага өөрчлөлт коэффициентүүдийг их хэмжээгээр өөрчилдөг.
\end{frame}

\end{document}
